\section{Why the labels satisfy every witness}
\label{sec:correctness}

Termination supplies a well-founded order for the correctness proof.
At a nonbase call, assume that all children satisfy the simultaneous
labelling guarantee. We will show first that the marks $I_-$ and
$I_+$ can safely receive their prescribed labels. We will then show
that every witness with a remaining claim survives the mass change.
These two statements allow the continuation to finish the call.

In the next three lemmas, fix a nonbase input $(A,D,a,b)$ and assume
Theorem~\ref{thm:labelling} for all its children. Witnesses, supports,
and claims without a prime refer to this parent input. Peaks and gaps
are those of Equations~\eqref{eq:max-gaps}--\eqref{eq:min-gaps}.

\subsection{Protection during motion and mass at stopping}

\begin{lemma}[The downward pass]
\label{lem:down}
For every relevant Max witness $x$, its peak $s_x(k)$ is constant
throughout Step~3 and at most $D/2$. In particular, no Max witness
claims an index in $I_-$. For every Min witness $y$ that has a claim
outside $I_-$,
\begin{equation}
             b\bigl(\{i:g_i^y(k^{\downarrow})<q\}\bigr)>3/4.
 \label{eq:down-heavy}
\end{equation}
\end{lemma}

\begin{proof}
Fix $x$ and consider one downward update from pivot $k$. Its
width-$q$ translation is a witness for the child, since that child
keeps Max masses $a$. By Lemma~\ref{lem:translation}, every
coordinate of gap at most $q/32=p$ is a Max claim of the translated
witness. The child's guarantee labels these coordinates $+$, and
they do not move. In particular at least one gap-zero coordinate
continues to attain the old peak.

At every other coordinate, the downward movement $\delta_i$ lies
in $[0,p]$. For its gap $g_i^x(k)>p$,
\[
 x_i-k_i^{\mathrm{new}}
           =s_x(k)-g_i^x(k)+\delta_i<s_x(k).
\]
Thus no coordinate can overtake the old peak, even though many may
move in the same update. The peak is unchanged. Initially $k=c$
and $x\le A+D$, so $s_x(c)\le D/2$. If $x$ claims $i$,
\[
 k_i^{\downarrow}\ge x_i-D/2
                    \ge A_i+D/2-P=c_i-P,
\]
which excludes $i$ from $I_-$.

For the mass assertion, let $y$ claim $j\notin I_-$. The stopped
pivot satisfies $k_j^{\downarrow}\ge c_j-P$, and the claim gives
$y_j\le A_j+P$. Hence
\begin{equation}
 s_y(k^{\downarrow})\ge D/2-2P>D/2-M.
 \label{eq:down-uncapped}
\end{equation}
A coordinate at the lower cap would have
$k_i^{\downarrow}-y_i\le D/2-M$, since $y_i\ge A_i$.
Every maximizer of $k_i^{\downarrow}-y_i$ is therefore strictly
above its cap.

Suppose the mass in Equation~\eqref{eq:down-heavy} were at most
$3/4$. The width-$q$ Min translation has exactly that support, so
it would have support mass at most one for the last child's masses
$\frac43 b$. Its gap-zero maximizers would all have to receive
$-$. At least one of those coordinates would then move down by a
positive amount, because it is above the cap and $p>0$. This
contradicts the final unchanged update. Equality at $3/4$ also
leads to this contradiction, proving the strict mass bound.
\end{proof}

The downward marks are now safe, but the pivot will still move upward.
The next lemma both protects the other type of claim and shows why the
first mass certificate survives that later movement.

\begin{lemma}[The upward pass]
\label{lem:up}
For every relevant Min witness $y$, its peak $s_y(k)$ is constant
throughout Step~4 and at most $D/2$. Thus no Min witness claims an
index in $I_+$. At the final pivot,
\begin{align}
 b\bigl(\{i:g_i^y(k^f)<q\}\bigr)&>3/4
   &&\text{if }y\text{ has a claim outside }I_-;
   \label{eq:final-min-heavy}\\
 a\bigl(\{i:g_i^x(k^f)<q\}\bigr)&>3/4
   &&\text{if }x\text{ has a claim outside }I_+.
   \label{eq:final-max-heavy}
\end{align}
\end{lemma}

\begin{proof}
In every upward child the Min masses are unchanged. Translate a fixed
relevant Min witness $y$ to that width-$q$ box. All gaps at most
$p$ become Min claims, so those coordinates are labelled $-$ and
stay fixed. Every other coordinate has gap greater than $p$ and
moves upward by at most $p$. Its new deviation $k_i^{\mathrm{new}}-y_i$
is therefore strictly below the old peak. A gap-zero coordinate does
not move, proving exact peak preservation. At the start,
$k^{\downarrow}\le c$ and $y\ge A$ imply $s_y\le D/2$.
A Min claim at $i$ consequently gives
\[
                    k_i^f\le y_i+D/2\le c_i+P,
\]
so it is outside $I_+$.

Each coordinate $k_i$ is nondecreasing during this pass while $s_y$
is fixed. It follows that
\begin{equation}
                  g_i^y(k^f)\le g_i^y(k^{\downarrow}).
 \label{eq:min-gap-persistence}
\end{equation}
The set of Min gaps $<q$ can only grow. Its $b$-mass was greater
than $3/4$ at $k^{\downarrow}$ whenever $y$ had a claim outside
$I_-$, by Lemma~\ref{lem:down}. Positivity of $b$ proves
Equation~\eqref{eq:final-min-heavy}.

It remains to obtain the Max mass certificate at this second stopping
point. If $x$ claims $j\notin I_+$, then
$x_j\ge A_j+D-P$ and $k_j^f\le c_j+P$, giving
\[
                      s_x(k^f)\ge D/2-2P>D/2-M.
\]
At an upper-capped coordinate, the inequalities
$k_i^f=c_i+M$ and $x_i\le A_i+D$ would give a deviation at
most $D/2-M$. Thus every maximizer lies strictly below the upper
cap. If the $a$-mass of the gaps $<q$ were at most $3/4$, the
width-$q$ translation would be a Max witness for the final child's
masses $\frac43 a$. Its maximizers would receive $+$ and force a
positive upward movement, contradicting the unchanged update.
This proves Equation~\eqref{eq:final-max-heavy}.
\end{proof}

The mass bounds are absolute: each qualifying set has mass greater
than $3/4$, while the entire witness support has mass at most one.
They are not merely fractional lower bounds relative to a support
that might have much smaller mass. This strength is exactly what
the reweighting calculation needs.

\subsection{Preserving the unsettled witnesses}

The provisional child does not directly label the parent's claims.
Its purpose is to assign the favorable label on each large-mass set
identified above. The larger width $D/2$ makes these sets claims of
the translated witnesses.

\begin{lemma}[Witness survival]
\label{lem:reweight}
Every parent Max witness with a claim outside $I_+$ remains a Max
witness for $(A,D,\bar a,\bar b)$. Every parent Min witness with
a claim outside $I_-$ remains a Min witness for that input.
\end{lemma}

\begin{proof}
Write $r=D/2$ for the provisional width. Since
\[
                         q=D/128<r/32=D/64,
\]
the translated width-$r$ witness claims every coordinate whose parent
gap is $<q$.

For a Max witness $x$ as in the statement, let
$N_x=\{i:g_i^x(k^f)<q\}$ and
$Q_+=\{i:\lambda_i^0=+\}$. Lemma~\ref{lem:up} gives
$a(N_x)>3/4$. By Lemma~\ref{lem:translation}, $N_x\subseteq S_x$;
its coordinates are claims of the width-$r$ translation, which
qualifies for the unchanged masses $a$. Correctness of that child
therefore implies $N_x\subseteq Q_+$. Reweighting the original
support gives
\begin{align*}
 \bar a(S_x)
   &=\frac85 a(S_x)-\frac45 a(S_x\cap Q_+)\\
   &<\frac85\cdot1-\frac45\cdot\frac34=1.
\end{align*}
The continuation has the same box and operator, so the subsolution
inequalities are unchanged. The displayed support bound is all that
is needed to keep $x$ a witness.

For Min set $N_y=\{i:g_i^y(k^f)<q\}$ and
$Q_-=\{i:\lambda_i^0=-\}$. The same two lemmas show
$N_y\subseteq S_y\cap Q_-$ and $b(N_y)>3/4$. Hence
\[
 \bar b(S_y)=\frac85 b(S_y)-\frac45 b(S_y\cap Q_-)
                         <\frac85-\frac45\cdot\frac34=1.
\]
The supersolution inequalities still hold in the unchanged box,
so $y$ also survives.
\end{proof}

\subsection{Completing the simultaneous induction}

\begin{proof}[Proof of Theorem~\ref{thm:labelling}]
Use induction on $R=d+B(\mu)$ from
Lemma~\ref{lem:termination}, which has already established
termination. We prove that a single returned vector meets every
parent obligation.

In the mass base case, $a_i b_i>1$ at every vertex. If $b_i>1$,
no Min witness can contain $i$ in its support, and hence no Min
witness can claim it; returning $+$ is permitted. If $b_i\le1$,
then $a_i>1$, so no Max witness can claim that vertex; returning
$-$ is permitted. This proves all simultaneous obligations in
Step~1.

In the width base case, comparison with the computed box solution
$z$ gives $x\le z\le y$ for every Max witness $x$ and Min
witness $y$. At a Max claim,
\[
                   z_i\ge A_i+D-P>A_i+D/2,
\]
and Step~2 returns $+$. At a Min claim,
$z_i\le A_i+P<A_i+D/2$, and it returns $-$. The argument
also applies to $D=1$, for which the midpoint is a rational
threshold rather than an integer coordinate.

In a nonbase call, all children have smaller rank and satisfy the
simultaneous guarantee by induction. The hypotheses of
Lemmas~\ref{lem:down}--\ref{lem:reweight} therefore hold. Fix
any Max witness and any index $i$ it claims. Lemma~\ref{lem:down}
excludes $i$ from $I_-$. If $i\in I_+$, the positive override
settles it. Otherwise that witness has a claim outside $I_+$ and
survives reweighting by Lemma~\ref{lem:reweight}. The continuation
must label $i$ positively, and no override changes it.

For any Min claim, Lemma~\ref{lem:up} excludes $I_+$. A claim
in $I_-$ receives the negative override. A claim outside $I_-$
belongs to a surviving Min witness and receives $-$ from the
continuation, with no later change. The marks $I_-$ refer to the
earlier pivot, but their safety was proved for the fixed parent
witnesses and remains valid after the upward pass. The two marked
sets may intersect; that intersection contains no claim of either
type, so the stipulated negative precedence causes no conflict.

Every vertex receives exactly one label. Since the witnesses and
claimed indices in the preceding argument were arbitrary, all their
obligations hold together. This completes the induction.
\end{proof}
